\documentclass[xcolor=dvipsnames,t,aspectratio=169]{beamer}

\usecolortheme{rose}
\usecolortheme{dolphin}
\usetheme{Boadilla}

\input{../../templates/slides/imports}
\input{../../templates/slides/settings}
\input{../../templates/slides/commands}

\titlegraphic{
    \includegraphics[scale = 0.5]{../../templates/slides/logo}
}

\logo{
\begin{tikzpicture}[overlay,remember picture]
\node[xshift=-.75cm, yshift=-1cm] at (current page.30){
    \includegraphics[width=0.1\textwidth]{../../templates/slides/logo}};
\end{tikzpicture}
}
\usetikzlibrary{positioning}

\newcommand{\highlight}[1]{{\color{nes_dark_orange} #1}}

\title{Pandas: manipulação e limpeza de dados}

\author{Eduardo Adame}

\date{{\color{nes_dark_purple}\textbf{Ciência de Dados}\\[0.5em] 27 de março de 2026 }}

\begin{document}

\frame[plain]{\titlepage}
\setcounter{framenumber}{0}


% ============================================================
\section{Introdução}
% ============================================================

\begin{frame}[c]{Por que Pandas?}

NumPy é excelente para álgebra — mas dados reais têm \highlight{colunas heterogêneas}, \highlight{rótulos}, \highlight{valores faltantes} e \highlight{estrutura tabular}.

% Sugestão de imagem: tabela CSV aberta no Excel / VS Code ao lado de um DataFrame renderizado no Jupyter — para ilustrar "dados no mundo real vs. dados no código"

\begin{itemize}
    \item Pandas adiciona \highlight{rótulos} a arrays NumPy
    \item Suporte nativo a dados faltantes (\texttt{NaN})
    \item Leitura/escrita de CSV, Excel, SQL, JSON, \ldots
    \item Operações de agrupamento, merge e reshape
\end{itemize}

\begin{display}[Regra de ouro]
NumPy para matrizes numéricas. Pandas para dados tabulares do mundo real.
\end{display}

\end{frame}


\begin{frame}[c]{As duas estruturas principais}

\begin{columns}
\column{0.48\textwidth}
\textbf{\color{nes_dark_purple}Series}
\begin{itemize}
    \item Array 1-D com \highlight{índice}
    \item Uma coluna de uma tabela
    \item \texttt{pd.Series}
\end{itemize}

\column{0.48\textwidth}
\textbf{\color{nes_dark_purple}DataFrame}
\begin{itemize}
    \item Tabela 2-D com \highlight{índice} e \highlight{colunas}
    \item Coleção de Series
    \item \texttt{pd.DataFrame}
\end{itemize}
\end{columns}

\vspace{1em}

% Sugestão de imagem: diagrama mostrando que um DataFrame é composto por várias Series verticais com um índice compartilhado — estilo "anatomia do DataFrame"

\begin{display}[Importante]
Todo DataFrame é um dicionário de Series com índice compartilhado.
\end{display}

\end{frame}


% ============================================================
\section{Series}
% ============================================================

\begin{frame}[c, fragile]{Criando uma Series}

\begin{code}[A partir de lista]{python}
import pandas as pd

s = pd.Series([10, 20, 30])
s = pd.Series([10, 20, 30], index=["a", "b", "c"])
\end{code}

\begin{code}[A partir de dicionário]{python}
s = pd.Series({"rio": 6.7, "sp": 12.3, "bh": 2.7})
\end{code}

\begin{itemize}
    \item O índice pode ser \highlight{inteiro, string ou datetime}
    \item \texttt{s.values} retorna o array NumPy subjacente
    \item \texttt{s.index} retorna o objeto índice
\end{itemize}

\end{frame}


\begin{frame}[c, fragile]{Indexação em Series}

\begin{code}[Por rótulo e por posição]{python}
s = pd.Series([10, 20, 30], index=["a", "b", "c"])

s["b"]          # por rótulo  -> 20
s.iloc[1]       # por posição -> 20
s[["a", "c"]]   # seleção múltipla
s["a":"b"]      # slice por rótulo (inclusivo)
\end{code}

\begin{display}[Atenção]
\texttt{s[1]} pode ser ambíguo. Prefira \texttt{.loc} (rótulo) ou \texttt{.iloc} (posição) explicitamente.
\end{display}

\end{frame}


\begin{frame}[c, fragile]{Operações e vetorização}

\begin{code}[Operações elementwise]{python}
s = pd.Series([1, 2, 3, 4])

s * 2
s ** 2
s + s           # alinhamento por índice!
\end{code}

\begin{code}[Funções de redução]{python}
s.sum()
s.mean()
s.describe()    # estatísticas resumidas
\end{code}

\begin{itemize}
    \item Pandas \highlight{alinha pelo índice} antes de operar
    \item Índices incompatíveis geram \texttt{NaN} automaticamente
\end{itemize}

\end{frame}


\begin{frame}[c, fragile]{Valores faltantes}

\begin{code}[Detectar e tratar]{python}
s = pd.Series([1.0, None, 3.0, None])

s.isna()            # máscara booleana
s.dropna()          # remove NaN
s.fillna(0)         # substitui NaN por 0
s.fillna(s.mean())  # substitui pela média
\end{code}

\begin{itemize}
    \item \texttt{NaN} é \highlight{contagioso}: operações com \texttt{NaN} retornam \texttt{NaN}
    \item \texttt{sum()}, \texttt{mean()} ignoram \texttt{NaN} por padrão
\end{itemize}

\begin{display}[Regra prática]
Sempre verifique \texttt{s.isna().sum()} antes de analisar dados novos.
\end{display}

\end{frame}


\begin{frame}[c, fragile]{\texttt{apply()} em Series}

\begin{code}[Aplicar função personalizada]{python}
s = pd.Series(["  Rio  ", "São Paulo", "belo horizonte"])

s.apply(str.strip)
s.apply(str.title)
s.apply(lambda x: len(x))
\end{code}

\begin{code}[Alternativa vetorizada para strings]{python}
s.str.strip()
s.str.title()
s.str.len()
\end{code}

\begin{itemize}
    \item \texttt{apply()} aceita qualquer função Python
    \item \highlight{Prefira} \texttt{.str}, \texttt{.dt} quando disponível — mais rápido
\end{itemize}

\end{frame}


% ============================================================
\section{DataFrame}
% ============================================================

\begin{frame}[c, fragile]{Criando um DataFrame}

\begin{code}[A partir de dicionário]{python}
df = pd.DataFrame({
    "nome":  ["Ana", "Bruno", "Carla"],
    "idade": [28, 34, 22],
    "nota":  [8.5, 7.0, 9.2],
})
\end{code}

\begin{code}[Leitura de CSV]{python}
df = pd.read_csv("dados.csv")
df.to_csv("saida.csv", index=False)
\end{code}

\begin{itemize}
    \item \texttt{df.shape}, \texttt{df.dtypes}, \texttt{df.head()}, \texttt{df.info()}
    \item Inspecione \highlight{sempre} antes de analisar
\end{itemize}

\end{frame}


\begin{frame}[c, fragile]{Indexação em DataFrame}

\begin{code}[Colunas e linhas]{python}
df["nome"]              # Series da coluna
df[["nome", "nota"]]    # sub-DataFrame

df.loc[0]               # linha por rótulo
df.loc[0, "nota"]       # célula por rótulo
df.iloc[0, 2]           # célula por posição
\end{code}

\begin{code}[Filtro booleano]{python}
df[df["nota"] > 8]
df[(df["idade"] < 30) & (df["nota"] > 8)]
\end{code}

\begin{display}[Atenção]
\texttt{df["col"]} retorna uma \texttt{Series}. \texttt{df[["col"]]} retorna um \texttt{DataFrame}.
\end{display}

\end{frame}


\begin{frame}[c, fragile]{Operações e \texttt{apply()}}

\begin{code}[Criar e transformar colunas]{python}
df["aprovado"] = df["nota"] >= 7.0
df["nota_norm"] = (df["nota"] - df["nota"].min()) \
                / (df["nota"].max() - df["nota"].min())
\end{code}

\begin{code}[apply() por coluna ou linha]{python}
df["nota"].apply(lambda x: "A" if x >= 9 else "B")

df.apply(lambda row: row["nota"] * 1.1, axis=1)
\end{code}

\begin{itemize}
    \item \texttt{axis=0} opera sobre linhas (padrão)
    \item \texttt{axis=1} opera sobre colunas da linha
\end{itemize}

\end{frame}


\begin{frame}[c, fragile]{\texttt{merge()} — combinando tabelas}

\begin{code}[Inner join]{python}
alunos   = pd.DataFrame({"id": [1,2,3], "nome": ["A","B","C"]})
notas    = pd.DataFrame({"id": [1,2,4], "nota": [8,7,9]})

pd.merge(alunos, notas, on="id")             # inner (padrão)
pd.merge(alunos, notas, on="id", how="left") # left join
\end{code}

\begin{itemize}
    \item \texttt{how}: \texttt{"inner"}, \texttt{"left"}, \texttt{"right"}, \texttt{"outer"}
    \item Equivalente ao JOIN do SQL
    \item Use \texttt{on} para chave comum; \texttt{left\_on}/\texttt{right\_on} para nomes distintos
\end{itemize}

\end{frame}


\begin{frame}[c, fragile]{Agregação e \texttt{groupby()}}

\begin{code}[Agrupamento básico]{python}
df.groupby("turma")["nota"].mean()
df.groupby("turma")["nota"].agg(["mean", "std", "count"])
\end{code}

\begin{code}[Múltiplas colunas]{python}
df.groupby(["turma", "aprovado"]).agg(
    media=("nota", "mean"),
    total=("nota", "count"),
)
\end{code}

\begin{display}[Padrão Split–Apply–Combine]
\texttt{groupby()} divide os dados, aplica uma função a cada grupo, e combina os resultados.
\end{display}

\end{frame}


% ============================================================
\section{Avançado}
% ============================================================

\begin{frame}[c, fragile]{Strings e datas}

\begin{code}[Accessor .str]{python}
df["nome"].str.lower()
df["nome"].str.contains("Silva")
df["nome"].str.replace("Dr. ", "", regex=False)
\end{code}

\begin{code}[Accessor .dt]{python}
df["data"] = pd.to_datetime(df["data"])

df["data"].dt.year
df["data"].dt.month
df["data"].dt.day_name()
\end{code}

\begin{itemize}
    \item \texttt{.str} e \texttt{.dt} são \highlight{vetorizados} — evite \texttt{apply()} para isso
\end{itemize}

\end{frame}


\begin{frame}[c, fragile]{Reshape: \texttt{pivot\_table()} e \texttt{melt()}}

\begin{code}[De longo para largo]{python}
df.pivot_table(
    values="nota",
    index="aluno",
    columns="disciplina",
    aggfunc="mean",
)
\end{code}

\begin{code}[De largo para longo]{python}
pd.melt(df, id_vars=["aluno"],
        value_vars=["mat", "por"],
        var_name="disciplina",
        value_name="nota")
\end{code}

% Sugestão de imagem: diagrama visual mostrando pivot (largo) e melt (longo) com tabelas de exemplo lado a lado — estilo "wide vs. long format"

\end{frame}


\begin{frame}[c, fragile]{View vs. Copy — armadilha clássica}

\begin{code}[Problema]{python}
sub = df[df["nota"] > 7]
sub["nota"] = 10        # SettingWithCopyWarning!
\end{code}

\begin{code}[Solução]{python}
sub = df[df["nota"] > 7].copy()
sub["nota"] = 10        # seguro
\end{code}

\begin{display}[Regra]
Sempre use \texttt{.copy()} ao criar um subconjunto que você pretende modificar.
\end{display}

\end{frame}


\begin{frame}[c, fragile]{Limpeza de dados — checklist}

\begin{code}[Inspeção inicial]{python}
df.info()           # tipos e não-nulos
df.isna().sum()     # contagem de NaN por coluna
df.duplicated().sum()
df.describe()
\end{code}

\begin{code}[Limpeza]{python}
df = df.drop_duplicates()
df = df.dropna(subset=["coluna_chave"])
df["coluna"] = df["coluna"].astype(float)
df.columns = df.columns.str.strip().str.lower()
\end{code}

\begin{itemize}
    \item Padronize nomes de colunas \highlight{primeiro}
    \item Corrija tipos \highlight{antes} de calcular qualquer coisa
\end{itemize}

\end{frame}


\begin{frame}[c, fragile]{Pipeline de análise}

\begin{code}[Encadeamento com method chaining]{python}
resultado = (
    pd.read_csv("vendas.csv")
      .rename(columns=str.lower)
      .dropna(subset=["preco", "quantidade"])
      .assign(total=lambda df: df["preco"] * df["quantidade"])
      .groupby("categoria")["total"]
      .sum()
      .sort_values(ascending=False)
)
\end{code}

\begin{itemize}
    \item \highlight{Method chaining} torna pipelines legíveis e reproduzíveis
    \item Use \texttt{.assign()} para criar colunas sem quebrar a cadeia
\end{itemize}

\end{frame}


\begin{frame}[c, noframenumbering, plain]
    \frametitle{~}
        \vfill
        \begin{center}
            {\Huge Obrigado!}\vspace{1.5em}\\
            {\Large \highlight{Alguma dúvida?}}\\
        \end{center}
        \vfill
        \begin{center}
            {\small Agora vamos para os exercícios!}
        \end{center}
\end{frame}

\end{document}